\chapter{2020年北京卷}

\section{单选题}

\begin{problem}
已知集合 \(A=\left\{ -1,0,1,2 \right\}\),\(B=\left\{ x\mid 0<x<3 \right\}\),则 \(A\cap B=\)(\quad).

\choices
    {\(\left\{ -1,0,1 \right\}\)}
    {\(\left\{ 0,1 \right\}\)}
    {\(\left\{ -1,1,2 \right\}\)}
    {\(\left\{ 1,2 \right\}\)}
\end{problem}

\begin{problem}
在复平面内,复数 \(z\) 对应的点的坐标是 \((1,2)\),则 \(\i z=\)(\quad).

\choices
    {\(1+2\i\)}
    {\(-2+\i\)}
    {\(1-2\i\)}
    {\(-2-\i\)}
\end{problem}

\begin{problem}
在 \(\left(\sqrt{x}-2\right)^5\) 的展开式中,\(x^2\) 的系数为(\quad).

\choices
    {\(-5\)}
    {\(5\)}
    {\(-10\)}
    {\(10\)}
\end{problem}

\begin{problem}
某三棱柱的底面为正三角形,其三视图如图所示,该三棱柱的表面积为(\quad).

\begin{center}
\bitmapfigure[max width=0.42\linewidth]{img/2020/beijing/q4_fig1.png}
\end{center}

\choices
    {\(6+3\sqrt3\)}
    {\(6+2\sqrt3\)}
    {\(12+3\sqrt3\)}
    {\(12+2\sqrt3\)}
\end{problem}

\begin{problem}
已知半径为 \(1\) 的圆经过点 \((3,4)\),则其圆心到原点的距离的最小值为(\quad).

\choices
    {\(4\)}
    {\(5\)}
    {\(6\)}
    {\(7\)}
\end{problem}

\begin{problem}
已知函数 \(f(x)=2^x-x-1\),则不等式 \(f(x)>0\) 的解集是(\quad).

\choices
    {\((-1,1)\)}
    {\((-\infty,-1)\cup(1,+\infty)\)}
    {\((0,1)\)}
    {\((-\infty,0)\cup(1,+\infty)\)}
\end{problem}

\begin{problem}
设抛物线的顶点为 \(O\),焦点为 \(F\),准线为 \(l\). \(P\) 是抛物线上异于 \(O\) 的一点,过 \(P\) 作 \(PQ\perp l\) 于 \(Q\),则线段 \(FQ\) 的垂直平分线(\quad).

\choices
    {经过点 \(O\)}
    {经过点 \(P\)}
    {平行于直线 \(OP\)}
    {垂直于直线 \(OP\)}
\end{problem}

\begin{problem}
在等差数列 \(\{a_n\}\) 中,\(a_1=-9\),\(a_3=-1\). 记 \(T_n=a_1a_2\cdots a_n\)(\(n=1\),\(2\),\(\dots\)),则数列 \(\{T_n\}\)(\quad).

\choices
    {有最大项,有最小项}
    {有最大项,无最小项}
    {无最大项,有最小项}
    {无最大项,无最小项}
\end{problem}

\begin{problem}
已知 \(\alpha\),\(\beta\in\R\),则``存在 \(k\in\Z\) 使得 \(\alpha=k\pi+(-1)^k\beta\)''是``\(\sin\alpha=\sin\beta\)''的(\quad).

\choices
    {充分而不必要条件}
    {必要而不充分条件}
    {充分必要条件}
    {既不充分也不必要条件}
\end{problem}

\begin{problem}
2020 年 3 月 14 日是全球首个国际圆周率日(\(\pi\) Day). 历史上,求圆周率 \(\pi\) 的方法有多种,与中国传统数学中的``割圆术''相似. 数学家阿尔·卡西的方法是:当正整数 \(n\) 充分大时,计算单位圆的内接正 \(6n\) 边形的周长和外切正 \(6n\) 边形(各边均与圆相切的正 \(6n\) 边形)的周长,将它们的算术平均数作为 \(2\pi\) 的近似值. 按照阿尔·卡西的方法,\(\pi\) 的近似值的表达式是(\quad).

\choices
    {\(3n\left(\sin\frac{30^\circ}{n}+\tan\frac{30^\circ}{n}\right)\)}
    {\(6n\left(\sin\frac{30^\circ}{n}+\tan\frac{30^\circ}{n}\right)\)}
    {\(3n\left(\sin\frac{60^\circ}{n}+\tan\frac{60^\circ}{n}\right)\)}
    {\(6n\left(\sin\frac{60^\circ}{n}+\tan\frac{60^\circ}{n}\right)\)}
\end{problem}

\section{填空题}

\begin{problem}
函数 \(f(x)=\dfrac{1}{x+1}+\ln x\) 的定义域是\fillinblank{}.
\end{problem}

\begin{problem}
已知双曲线 \(C:\dfrac{x^2}{6}-\dfrac{y^2}{3}=1\),则 \(C\) 的右焦点的坐标为\fillinblank{};\(C\) 的焦点到其渐近线的距离是\fillinblank{}.
\end{problem}

\begin{problem}
已知正方形 \(ABCD\) 的边长为 \(2\),点 \(P\) 满足 \(\overrightarrow{AP}=\dfrac12\left(\overrightarrow{AB}+\overrightarrow{AC}\right)\),则 \(|PD|=\fillinblank{}\);\(\overrightarrow{PB}\cdot\overrightarrow{PD}=\fillinblank{}\).
\end{problem}

\begin{problem}
若函数 \(f(x)=\sin(x+\varphi)+\cos x\) 的最大值为 \(2\),则常数 \(\varphi\) 的一个取值为\fillinblank{}.
\end{problem}

\begin{problem}
为满足人民对美好生活的向往,环保部门要求相关企业加强污水治理,排放未达标的企业要限期整改,设企业的污水排放量 \(W\) 与时间 \(t\) 的关系为 \(W=f(t)\),用 \(-\dfrac{f(b)-f(a)}{b-a}\) 的大小评价在 \([a,b]\) 这段时间内企业污水治理能力的强弱,已知整改期内,甲、乙两企业的污水排放量与时间的关系如下图所示.

给出下列四个结论:

\begin{circlelist}
\item 在 \([t_1,t_2]\) 这段时间内,甲企业的污水治理能力比乙企业强;

\item 在 \(t_2\) 时刻,甲企业的污水治理能力比乙企业强;

\item 在 \(t_3\) 时刻,甲、乙两企业的污水排放都已达标;

\item 甲企业在 \([0,t_1]\),\([t_1,t_2]\),\([t_2,t_3]\) 这三段时间中,在 \([0,t_1]\) 的污水治理能力最强.
\end{circlelist}
其中所有正确结论的序号是\fillinblank{}.

\begin{center}
\FigureLayoutDeclare{img/2020/beijing/q15_fig1.png}{17.0cm}{78bp 29bp 117bp 7bp}
\bitmapfigure[max width=0.80\linewidth]{img/2020/beijing/q15_fig1.png}
\end{center}
\end{problem}

\section{解答题}

\begin{problem}
如图,在正方体 \(ABCD-A_1B_1C_1D_1\) 中,\(E\) 为 \(BB_1\) 的中点.

\begin{enumerate}
    \item 求证:\(BC_1\parallel\) 平面 \(AD_1E\);
    \item 求直线 \(AA_1\) 与平面 \(AD_1E\) 所成角的正弦值.
\end{enumerate}

\begin{center}
\bitmapfigure[max width=0.60\linewidth]{img/2020/beijing/q16_fig1.png}
\end{center}
\end{problem}

\begin{problem}
在 \(\triangle ABC\) 中,\(a+b=11\),再从条件\(\circled{1}\)、条件\(\circled{2}\)这两个条件中选择一个作为已知,求:

\begin{enumerate}
    \item \(a\) 的值;
    \item \(\sin C\) 和 \(\triangle ABC\) 的面积.
\end{enumerate}

\begin{circlelist}
\item \(c=7\),\(\cos A=-\dfrac17\);

\item \(\cos A=\dfrac18\),\(\cos B=\dfrac{9}{16}\).
\end{circlelist}
\end{problem}

\begin{problem}
某校为举办甲、乙两项不同活动,分别设计了相应的活动方案:方案一、方案二. 为了解该校学生对活动方案是否支持,对学生进行简单随机抽样,获得数据如下表:

\begin{center}
\begin{tabular}{c|cc|cc}
\hline
& \multicolumn{2}{c|}{男生} & \multicolumn{2}{c}{女生} \\
\hline
& 支持 & 不支持 & 支持 & 不支持 \\
\hline
方案一 & \(200\text{ 人}\) & \(400\text{ 人}\) & \(300\text{ 人}\) & \(100\text{ 人}\) \\
方案二 & \(350\text{ 人}\) & \(250\text{ 人}\) & \(150\text{ 人}\) & \(250\text{ 人}\) \\
\hline
\end{tabular}
\end{center}

假设所有学生对活动方案是否支持相互独立.

\begin{enumerate}
    \item 分别估计该校男生支持方案一的概率,该校女生支持方案一的概率;
    \item 从该校全体男生中随机抽取 \(2\text{ 人}\),全体女生中随机抽取 \(1\text{ 人}\),估计这 \(3\text{ 人}\)中恰有 \(2\text{ 人}\)支持方案一的概率;
    \item 将该校学生支持方案的概率估计值记为 \(p_0\),假设该校年级有 \(500\text{ 名}\) 男生和 \(300\text{ 名}\) 女生,除一年级外其他年级学生支持方案二的概率估计值记为 \(p_1\),试比较 \(p_0\) 与 \(p_1\) 的大小. (结论不要求证明)
\end{enumerate}
\end{problem}

\begin{problem}
已知函数 \(f(x)=12-x^2\).

\begin{enumerate}
    \item 求曲线 \(y=f(x)\) 的斜率等于 \(-2\) 的切线方程;
    \item 设曲线 \(y=f(x)\) 在点 \((t,f(t))\) 处的切线与坐标轴围成的三角形的面积为 \(S(t)\),求 \(S(t)\) 的最小值.
\end{enumerate}
\end{problem}

\begin{problem}
已知椭圆 \(C:\dfrac{x^2}{a^2}+\dfrac{y^2}{b^2}=1\) 过点 \(A(-2,-1)\),且 \(a=2b\).

\begin{enumerate}
    \item 求椭圆 \(C\) 的方程;
    \item 过点 \(B(-4,0)\) 的直线 \(l\) 交椭圆 \(C\) 于点 \(M\),\(N\),直线 \(MA\),\(NA\) 分别交直线 \(x=-4\) 于点 \(P\),\(Q\). 求 \(\left|\dfrac{PB}{BQ}\right|\) 的值.
\end{enumerate}
\end{problem}

\begin{problem}
已知 \(\{a_n\}\) 是无穷数列. 给出两个性质:

\begin{circlelist}
\item 对于 \(\{a_n\}\) 中任意两项 \(a_i\),\(a_j\)(\(i>j\)),在 \(\{a_n\}\) 中都存在一项 \(a_m\),使得 \(\dfrac{a_i^2}{a_j}=a_m\);

\item 对于 \(\{a_n\}\) 中任意项 \(a_n\)(\(n\ge3\)),在 \(\{a_n\}\) 中都存在两项 \(a_k\),\(a_l\)(\(k>l\)),使得 \(a_n=\dfrac{a_k^2}{a_l}\).
\end{circlelist}

\begin{enumerate}
\item 若 \(a_n=n\)(\(n=1\),\(2\),\(\dots\)),判断数列 \(\{a_n\}\) 是否满足性质(1),说明理由;
\item 若 \(a_n=2^{n-1}\)(\(n=1\),\(2\),\(\dots\)),判断数列 \(\{a_n\}\) 是否同时满足性质(1)和性质(2),说明理由;
\item 若 \(\{a_n\}\) 是递增数列,且同时满足性质(1)和性质(2),证明:\(\{a_n\}\) 为等比数列.
\end{enumerate}
\end{problem}

